% Job posting: https://ca.linkedin.com/jobs/view/computer-vision-engineer-markham-on-at-tenth-revolution-group-4466906192
\documentclass[11pt]{article}
\usepackage[left=1in,top=1in,right=1in,bottom=1in]{geometry}
\usepackage[hidelinks]{hyperref}
\usepackage{setspace}
\usepackage[T1]{fontenc}
\usepackage{newtxtext,newtxmath}
\ifdefined\pdfgentounicode
  \input{glyphtounicode}
  \pdfgentounicode=1
\fi
\setstretch{1.1}
\pagenumbering{gobble}
\begin{document}
\pagestyle{empty}

\noindent Nishal Stanislaus Silva, PhD \\
Toronto, ON \\
nishal.silva@hotmail.com \,|\, +1 613 200 2030 \\
nishal.xyz \,|\, linkedin.com/in/nishal-silva \,|\, github.com/ns2max

\vspace{1.5em}

\noindent Dear Hiring Manager,

\vspace{1em}

\noindent I am writing to apply for the Computer Vision Engineer role in Markham that your firm is recruiting for. For 5.5 years at MAS Holdings, South Asia's largest apparel manufacturer, I built and ran computer-vision systems on live production floors, not in a lab: an automated multilingual care-label QC system using OpenCV template and feature matching with an explainable defect overlay, iterated from a handheld scanner to an industrial camera on a conveyor with PLC-controlled reject, cutting inspection time 99.5\% while keeping a human-in-the-loop path for borderline cases. This is the same discipline the posting describes: computer vision and pattern-recognition R\&D that has to survive contact with an automation line, not just a benchmark.

\vspace{1em}

\noindent Against the specific requirements: I have reverse-engineered PLC trigger points and coordinated closed-loop electrical/software subsystems in custom C++ for a direct-to-garment printer integration, and separately designed a geometry-alignment algorithm and built the RIP integration for Promptly, MAS's on-demand printing product now running in four countries, so I am comfortable owning the boundary between vision software and factory hardware. I also built a loom-side fabric-defect detection system using a microscopic camera array that cut operator headcount 50\%, and OCR/layout-segmentation pipelines for garment tech-pack digitization. My production C++ work has been on Linux and embedded toolchains rather than Visual Studio on Windows specifically, and I have not done PCB schematic capture in AutoCAD; both are a natural extension of the electrical/software integration work I already do, and I would expect to pick up the Windows-side toolchain quickly. I have also delivered hands-on CV training to engineering teams and represented my group at international trade shows, so customer-facing installation and commissioning work is familiar territory, including the travel that comes with it.

\vspace{1em}

\noindent Manufacturing-floor automation is where I have chosen to build my career, from apparel production in Colombo to a PhD adding formal rigor to how I benchmark and validate a system before it goes live. I would bring that same prototype-to-production discipline to your client's automation and computer-vision programs from day one.

\vspace{1em}

\noindent I am a Canadian Permanent Resident and do not require sponsorship. I am based in Toronto, a short commute from Markham, and available immediately, including for the periodic travel the role requires. I would welcome the chance to discuss how my industrial computer-vision background can contribute to your client's team.

\vspace{1.5em}

\noindent Sincerely, \\
Nishal Stanislaus Silva

\end{document}
